% !PRINTED 298
\weberpage{298}
Ebenso ist zu schliessen, dass $\psi$ in jedem der Intervalle
$1, 3, 5 \dots 4n - 1$ einmal und nur einmal durch Null geht.
Wir haben also eine Eintheilung der Kreisperipherie, wie sie im
vorigen Paragraphen zur Bestimmung der Charakteristik verlangt
wurde, und da in den Theilpunkten $\omega, 5\omega, 9\omega \dots (4n - 3)\omega$
das Product $\varphi\psi$ positiv ist, so kommen auf der Kreisperipherie
nur Punkte $E(f; \varphi, \psi)$ vor und ihre Anzahl ist $2n$; die Anzahl
der Wurzeln von $F(z)$ ist also gleich $n$, was zu beweisen war.
